%! Parametrically Defined Circles and Lines — Unit 4, Lesson 4.4 — Exit Ticket (Answer Key)
\documentclass[10pt]{article}
\usepackage{precalculus-article}
\usepackage{precalculus-key}   % pulls in -boxes; do NOT also load -boxes
\newcommand{\TallMath}[1]{$\displaystyle #1\rule[-1.4em]{0pt}{3.2em}$}

\begin{document}

\pageheader{Unit 4, Lesson 4.4}{Exit Ticket --- Answer Key}
\namedateperiod

\vspace{0.12in}

\begin{notesbox}{Parametrically Defined Circles and Lines --- Key}

\textbf{1.}\quad A circle is centered at $(3,\,-1)$ with radius $4$.
Write parametric equations for \emph{counterclockwise} motion that
starts at $(7,\,-1)$.

\vspace{4pt}
$x(t) = $ \ans{$3 + 4\cos t$} \qquad $y(t) = $ \ans{$-1 + 4\sin t$}

\vspace{0.04in}
Domain: \ans{$0 \le t \le 2\pi$}

\vspace{0.22in}
\textbf{2.}\quad Write a parametrization for the line segment from
$(0,\,6)$ to $(4,\,0)$.

\vspace{4pt}
$x(t) = $ \ans{$0 + 4t = 4t$} \qquad $y(t) = $ \ans{$6 + (0-6)t = 6 - 6t$}

\vspace{0.04in}
Domain: \ans{$0 \le t \le 1$}

\vspace{0.16in}
At what value of $t$ does the segment cross the $x$-axis (where $y=0$)?

\vspace{4pt}
Set $y(t) = 0$: $6 - 6t = 0 \;\Rightarrow\; 6t = 6 \;\Rightarrow\; t = $ \ans{$1$}.

\ans{The segment reaches the $x$-axis at $t = 1$, which is the endpoint $(4,\;0)$.}

\begin{teachernote}
  Q1 checks the transformation template directly. Q2 is purposely designed
  so the $x$-axis crossing is exactly the endpoint ($t=1$). If students
  get $t=1$, ask: ``Is that in the domain? What point does it correspond
  to?'' This builds the habit of checking domain validity.
\end{teachernote}

\end{notesbox}

\end{document}
